\documentclass[12pt]{article}
\usepackage[utf8]{inputenc}
\usepackage[spanish]{babel}
\usepackage{amsmath}
\usepackage{amssymb}
\usepackage{booktabs}
\usepackage{geometry}
\geometry{margin=2.5cm}

\title{Resumen: Cap\'itulo 4 --- Modelos con Tendencia\\[4pt]
\large Enders, \textit{Applied Econometric Time Series}, 4.a~ed.}
\author{}
\date{}

\begin{document}
\maketitle
\tableofcontents
\newpage

%%% =========================================================
\section{Tendencias Determin\'isticas y Estoc\'asticas}
%%% =========================================================

Una serie temporal puede exhibir dos tipos distintos de tendencia:

\textbf{Trend stationary (TS) / estacionario en tendencia:} proceso cuya
tendencia es una funci\'on determin\'istica del tiempo; sustrayendo dicha
tendencia se obtiene un proceso estacionario:

\begin{equation}
  y_t = y_0 + a_0\,t + A(L)\varepsilon_t
\end{equation}

donde $A(L)$ es un polinomio de retardos estacionario e invertible y
$\{\varepsilon_t\}$ es ruido blanco.

\textbf{Random walk / paseo aleatorio:} proceso con ra\'iz unitaria cuya
varianza crece linealmente con el tiempo:

\begin{equation}
  y_t = y_{t-1} + \varepsilon_t, \qquad \mathrm{Var}(y_t) = t\sigma^2
\end{equation}

La funci\'on de autocorrelaci\'on (ACF) satisface
$\rho_s \approx [(t-s)/t]^{1/2}$ y decae muy lentamente en muestras
finitas, pudiendo confundirse con la de un proceso estacionario de
memoria larga.

\textbf{Random walk plus drift / paseo aleatorio con deriva:}

\begin{equation}
  y_t = y_0 + a_0\,t + \sum_{i=1}^{t}\varepsilon_i
\end{equation}

La serie combina una tendencia determin\'istica $a_0 t$ con una tendencia
estoc\'astica acumulada $\sum \varepsilon_i$; visualmente puede confundirse
con un proceso TS.

\textbf{Random walk plus noise / paseo aleatorio m\'as ruido:} al a\~nadir
un componente de ruido $\eta_t$ al proceso integrado, las primeras
diferencias siguen un proceso ARIMA$(0,1,1)$; la presencia de un t\'ermino
de media m\'ovil negativo en $\Delta y_t$ es indicio de este tipo de proceso.

La forma general con componentes irregulares de estructura ARMA es:

\begin{equation}
  y_t = y_0 + a_0\,t + \sum_{i=1}^{t}\varepsilon_i + A(L)\eta_t
\end{equation}

\textbf{Unobserved components models / modelos de componentes no
observados:} marco en el que la tendencia, el ciclo y la componente
irregular son variables latentes. La inferencia sobre estos componentes
se denomina \textbf{signal extraction / extracci\'on de se\~nal}.

%%% =========================================================
\section{Eliminaci\'on de la Tendencia}
%%% =========================================================

El tratamiento correcto depende de la naturaleza de la tendencia:

\begin{itemize}
  \item \textbf{Differencing / diferenciaci\'on:} aplica el operador
        $\Delta = 1-L$ para eliminar tendencias estoc\'asticas. Diferenciar
        un proceso TS lo sobrediferencia e introduce ra\'ices MA no invertibles.
  \item \textbf{Detrending / eliminaci\'on de tendencia determin\'istica:}
        regresa la serie sobre funciones del tiempo. Desestacionalizar un
        proceso DS deja una tendencia estoc\'astica residual.
\end{itemize}

\textbf{Integrated of order $d$ / integrado de orden $d$, $I(d)$:} la
serie $\{y_t\}$ es $I(d)$ si su $d$-\'esima diferencia es estacionaria e
invertible; posee exactamente $d$ ra\'ices unitarias. Un proceso
ARIMA$(p,d,q)$ es $I(d)$.

\textbf{Difference stationary (DS) / estacionario en diferencias:}
proceso $I(1)$ que se estacionaliza mediante diferenciaci\'on.

Nelson y Plosser (1982) aplicaron pruebas de ra\'iz unitaria a 14 series
macroecon\'omicas de Estados Unidos y concluyeron que la mayor\'ia son DS,
en contraste con la visi\'on tradicional de series TS.

%%% =========================================================
\section{Ra\'ices Unitarias y Residuos de Regresi\'on}
%%% =========================================================

Considere la regresi\'on simple de $\{y_t\}$ sobre $\{z_t\}$:

\begin{equation}
  y_t = a_0 + a_1 z_t + e_t
\end{equation}

Granger y Newbold (1974) identificaron cuatro casos seg\'un el orden de
integraci\'on de ambas series:

\begin{enumerate}
  \item \textbf{Caso 1}: ambas estacionarias $\Rightarrow$ MCO v\'alido;
        $t$-estad\'isticos asint\'oticamente normales.
  \item \textbf{Caso 2}: \'ordenes distintos $\Rightarrow$ regresi\'on
        sin significado econ\'omico.
  \item \textbf{Caso 3}: mismo orden, residuo $\{e_t\}$ no estacionario
        $\Rightarrow$ \textbf{spurious regression / regresi\'on espuria}:
        $R^2$ elevado y $t$-estad\'isticos significativos sin relaci\'on
        econ\'omica genuina; se recomienda trabajar con primeras diferencias.
  \item \textbf{Caso 4}: mismo orden, residuo $\{e_t\}$ estacionario
        $\Rightarrow$ series \textbf{cointegrated / cointegradas}; MCO
        v\'alido con estimadores superconsistentes.
\end{enumerate}

%%% =========================================================
\section{El M\'etodo de Monte Carlo}
%%% =========================================================

\textbf{Monte Carlo experiment / experimento de Monte Carlo:} t\'ecnica de
simulaci\'on que genera $N$ muestras artificiales a partir de un
\textbf{data-generating process (DGP) / proceso generador de datos}
especificado y calcula el estad\'istico de inter\'es en cada r\'eplica.
La distribuci\'on emp\'irica resultante sirve como aproximaci\'on de la
distribuci\'on te\'orica del estad\'istico, cuya derivaci\'on anal\'itica
suele ser intratable.

La Ley de los Grandes N\'umeros justifica que, para $N$ suficientemente
grande, la distribuci\'on emp\'irica converge a la poblacional. Dickey y
Fuller emplearon este m\'etodo para calcular los valores cr\'iticos de sus
pruebas bajo la hip\'otesis nula de paseo aleatorio.

%%% =========================================================
\section{Pruebas de Dickey-Fuller}
%%% =========================================================

Las \textbf{Dickey-Fuller (DF) tests / pruebas de Dickey-Fuller} contrastan
$H_0\colon \gamma = 0$ (ra\'iz unitaria) frente a $H_1\colon \gamma < 0$
(estacionariedad). Bajo $H_0$, los estad\'isticos no siguen distribuciones
$t$ est\'andar; sus distribuciones asint\'oticas, obtenidas por simulaci\'on,
tienen valores cr\'iticos m\'as negativos que los de la normal.

Se distinguen tres regresiones auxiliares seg\'un los t\'erminos
determin\'isticos incluidos:

\textbf{Forma 1} --- paseo aleatorio puro (estad\'istico $\tau$):
\begin{equation}
  \Delta y_t = \gamma\,y_{t-1} + \varepsilon_t
\end{equation}

\textbf{Forma 2} --- con intercepto (estad\'istico $\tau_\mu$):
\begin{equation}
  \Delta y_t = a_0 + \gamma\,y_{t-1} + \varepsilon_t
\end{equation}

\textbf{Forma 3} --- con intercepto y tendencia lineal
(estad\'istico $\tau_\tau$):
\begin{equation}
  \Delta y_t = a_0 + \gamma\,y_{t-1} + a_2\,t + \varepsilon_t
\end{equation}

Se definen tambi\'en estad\'isticos $F$ conjuntos: $\phi_1$ (intercepto y
ra\'iz unitaria), $\phi_2$ (intercepto, tendencia y ra\'iz unitaria) y
$\phi_3$ (tendencia y ra\'iz unitaria), que permiten contrastar
hip\'otesis conjuntas sobre los par\'ametros determin\'isticos.

%%% =========================================================
\section{Aplicaciones de la Prueba de Dickey-Fuller}
%%% =========================================================

\subsection{Series de Nelson y Plosser (1982)}

Nelson y Plosser aplicaron las tres formas de la prueba DF a 14 series
macroecon\'omicas de Estados Unidos. No pudieron rechazar la hip\'otesis
nula de ra\'iz unitaria en 13 de las 14 series, apoyando la
caracterizaci\'on de las series macroecon\'omicas como procesos DS.

\subsection{Paridad de Poder de Compra}

La \textbf{Purchasing Power Parity (PPP) / Paridad de Poder de Compra}
postula que el tipo de cambio nominal refleja los diferenciales de precios
relativos:

\begin{equation}
  e_t = p_t - p_t^* + d_t
\end{equation}

donde $e_t$ es el logaritmo del tipo de cambio nominal, $p_t$ y $p_t^*$
son los logaritmos del nivel de precios dom\'estico y extranjero, y $d_t$
captura desviaciones transitorias. El \textbf{real exchange rate / tipo de
cambio real} se define como:

\begin{equation}
  r_t \;\equiv\; e_t + p_t^* - p_t
\end{equation}

La PPP se cumple a largo plazo si y solo si $\{r_t\}$ es estacionario.
La evidencia DF habitualmente no rechaza la ra\'iz unitaria en $\{r_t\}$,
aunque esto puede reflejar baja potencia m\'as que ausencia de PPP.

%%% =========================================================
\section{Extensiones de la Prueba de Dickey-Fuller}
%%% =========================================================

\subsection{Prueba Aumentada de Dickey-Fuller (ADF)}

La \textbf{Augmented Dickey-Fuller (ADF) test / prueba aumentada de
Dickey-Fuller} agrega retardos de $\Delta y_t$ para controlar la
correlaci\'on serial de los residuos:

\begin{equation}
  \Delta y_t = a_0 + \gamma\,y_{t-1} + a_2\,t
               + \sum_{i=1}^{p}\beta_i\,\Delta y_{t-i} + \varepsilon_t
\end{equation}

La selecci\'on de $p$ sigue la \textbf{general-to-specific methodology /
metodolog\'ia de lo general a lo espec\'ifico}: partir del n\'umero m\'aximo
de retardos y eliminar los no significativos secuencialmente.
Ng y Perron (2001) proponen el \textbf{Modified Akaike Information
Criterion (MAIC)}:

\begin{equation}
  \mathrm{MAIC} = T\ln\!\bigl(\widehat{\mathrm{SSR}}\bigr)
                  + 2n + 2\hat{\tau}(n)
\end{equation}

donde $\hat{\tau}(n)$ penaliza en funci\'on del coeficiente estimado de
$y_{t-1}$, mitigando la tendencia al sobrerechazo cuando el proceso
verdadero tiene t\'erminos MA negativos.

Said y Dickey (1984) demuestran que un proceso ARIMA$(p,1,q)$ general
puede aproximarse mediante un ARIMA$(n,1,0)$ con $n \leq T^{1/3}$,
validando el uso de la prueba ADF en contextos m\'as generales.

\subsection{Ra\'ices Unitarias M\'ultiples}

Dickey y Pantula (1987) proponen un procedimiento secuencial para detectar
\textbf{multiple unit roots / ra\'ices unitarias m\'ultiples}: probar
primero la hip\'otesis con mayor n\'umero de ra\'ices y avanzar hacia
hip\'otesis con menor n\'umero solo si se rechaza la anterior.

\subsection{Ra\'ices Unitarias Estacionales}

Las series estacionales pueden presentar \textbf{seasonal unit roots /
ra\'ices unitarias estacionales} a frecuencias distintas de cero.
Hylleberg, Engle, Granger y Yoo (HEGY, 1990) dise\~naron una prueba que
detecta ra\'ices unitarias en tres frecuencias: no estacional
($\gamma_1 = 0$), semestral ($\gamma_2 = 0$) y anual
($\gamma_5 = \gamma_6 = 0$ conjuntamente).

\subsection{T\'erminos MA Negativos y Sobrerechazo}

Cuando el proceso verdadero contiene t\'erminos MA negativos y se aproxima
mediante retardos AR, la prueba ADF tiende a sobrerechazar $H_0$. El
criterio MAIC corrige este sesgo seleccionando un mayor n\'umero de
retardos en presencia de t\'erminos MA negativos significativos.

%%% =========================================================
\section{Cambio Estructural}
%%% =========================================================

Los \textbf{structural breaks / quiebres estructurales} sesgan las pruebas
DF hacia la no rejecci\'on de la ra\'iz unitaria: un cambio en el nivel o
la pendiente de la tendencia produce ACF lentamente decrecientes incluso
en series estacionarias.

\textbf{Level dummy ($D_L$) / variable indicadora de nivel:} toma valor 0
antes del quiebre y 1 despu\'es; captura un cambio permanente en el nivel
medio de la serie.

\textbf{Pulse dummy ($D_P$) / variable indicadora de pulso:} toma valor 1
solo en el per\'iodo del quiebre y 0 en los dem\'as; captura un efecto
transitorio.

\subsection{Prueba de Perron (1989)}

Perron (1989) propone una extensi\'on formal que incorpora un quiebre
estructural \emph{conocido} en la regresi\'on auxiliar. El par\'ametro
$\lambda = \tau/T$ denota la proporci\'on de observaciones anteriores al
quiebre. Los valores cr\'iticos dependen de $\lambda$ y son m\'as
negativos que los de la prueba DF est\'andar; para $\lambda = 0.5$, el
valor cr\'itico al 5\,\% es aproximadamente $-3.76$, frente a $-3.41$
de la prueba DF sin quiebre.

Aplicando esta prueba a los datos de Nelson y Plosser, Perron concluy\'o
que la mayor\'ia de las series macroecon\'omicas son en realidad TS con
quiebres estructurales, revirtiendo la conclusi\'on original.

%%% =========================================================
\section{Potencia y los Regresores Determin\'isticos}
%%% =========================================================

\textbf{Power / potencia} de una prueba: probabilidad de rechazar
correctamente la hip\'otesis nula cuando esta es falsa; igual a
$1 - P(\text{error de tipo II})$. Las pruebas DF tienen potencia muy baja
cuando la ra\'iz autorregresiva es pr\'oxima, pero no igual, a la unidad
(\textbf{near unit root / ra\'iz casi unitaria}).

Campbell y Perron (1991) identifican dos fuentes de p\'erdida de potencia
vinculadas a los regresores determin\'isticos:

\begin{itemize}
  \item \textbf{Subincluci\'on:} omitir regresores determin\'isticos
        presentes en el DGP sesga el estimador de $\gamma$ y reduce la
        potencia.
  \item \textbf{Sobreincluci\'on:} incluir regresores determin\'isticos
        irrelevantes consume grados de libertad y desplaza los valores
        cr\'iticos, reduciendo tambi\'en la potencia.
\end{itemize}

El principio operativo es: las pruebas de ra\'iz unitaria se realizan
\emph{condicionales} a la especificaci\'on de los regresores
determin\'isticos, y las pruebas sobre dichos regresores se realizan
\emph{condicionales} a la presencia o ausencia de ra\'iz unitaria.

%%% =========================================================
\section{Pruebas con Mayor Potencia}
%%% =========================================================

\subsection{Prueba LM de Schmidt y Phillips (1992)}

Schmidt y Phillips (1992) proponen un procedimiento en dos etapas basado
en el principio del \textbf{Lagrange Multiplier (LM) / multiplicador de
Lagrange}. En la primera etapa se estima la tendencia determin\'istica
mediante la regresi\'on de $\Delta y_t$ sobre una constante, sin incluir
$y_{t-1}$, lo que evita contaminar la estimaci\'on de la tendencia con
el componente persistente. En la segunda etapa se realiza la prueba de
ra\'iz unitaria sobre los datos desestacionalizados. Al separar ambas
etapas, la prueba gana potencia respecto de la DF cl\'asica.

\subsection{Prueba ERS y DF-GLS (Elliott, Rothenberg y Stock, 1996)}

Elliott, Rothenberg y Stock (1996) proponen la prueba
\textbf{DF-GLS / Dickey-Fuller con m\'inimos cuadrados generalizados},
tambi\'en llamada prueba ERS. El procedimiento utiliza una
\textbf{near-differencing / cuasi-diferenciaci\'on} con par\'ametro
$\alpha$ pr\'oximo a la unidad:

\begin{itemize}
  \item Solo intercepto en el DGP: $\alpha = 1 - 7/T$.
  \item Intercepto y tendencia en el DGP: $\alpha = 1 - 13.5/T$.
\end{itemize}

La variable cuasi-diferenciada $\tilde{y}_t = y_t - \alpha\,y_{t-1}$ y
los regresores determin\'isticos cuasi-diferenciados se utilizan para
estimar de forma m\'as eficiente los par\'ametros de la tendencia. La
prueba de ra\'iz unitaria se aplica entonces a la serie desestacionalizada
mediante las mismas regresiones DF.

La prueba ERS tiene una potencia considerablemente mayor que la prueba DF
est\'andar, especialmente cuando la ra\'iz autorregresiva es pr\'oxima
a la unidad. Elliott (1999) propone una variante en la que la condici\'on
inicial se extrae de su distribuci\'on incondicional, lo que modifica
los valores cr\'iticos.

%%% =========================================================
\section{Pruebas de Ra\'iz Unitaria en Panel}
%%% =========================================================

Cuando se dispone de $n$ series similares y la potencia individual es
baja, es posible combinar la informaci\'on mediante una
\textbf{panel unit root test / prueba de ra\'iz unitaria en panel}.

\subsection{Prueba Im-Pesaran-Shin (IPS, 2002)}

Im, Pesaran y Shin (2002) proponen realizar una prueba ADF independiente
para cada serie $i = 1,\ldots,n$:

\begin{equation}
  \Delta y_{it} = a_{i0} + \gamma_i\,y_{i,t-1} + a_{i2}\,t
                  + \sum_{j=1}^{p_i}\beta_{ij}\,\Delta y_{it-j}
                  + \varepsilon_{it}
\end{equation}

Se obtiene el estad\'istico $t_i$ para $H_0\colon \gamma_i = 0$ en
cada ecuaci\'on y se forma el promedio:

\begin{equation}
  \bar{t} = \frac{1}{n}\sum_{i=1}^{n} t_i
\end{equation}

El estad\'istico estandarizado:

\begin{equation}
  Z_{\bar{t}} = \frac{\sqrt{n}\bigl[\bar{t} - E(\bar{t})\bigr]}
                     {\sqrt{\mathrm{Var}(\bar{t})}}
\end{equation}

converge asint\'oticamente a $\mathcal{N}(0,1)$ para $n$ y $T$ grandes.
Los valores $E(\bar{t})$ y $\mathrm{Var}(\bar{t})$ corrigen el sesgo a
la baja del estimador MCO de $\gamma_i$ y se tabulan por simulaci\'on.

La hip\'otesis nula es $\gamma_1 = \cdots = \gamma_n = 0$ (todas las series
son no estacionarias); el rechazo implica que \emph{al menos una} serie
es estacionaria, sin identificar cu\'ales.

\subsection{Correcci\'on por Efectos Temporales Comunes}

Si los residuos de las ecuaciones individuales est\'an correlacionados
contempor\'aneamente ($E\varepsilon_{it}\varepsilon_{jt} \neq 0$), los
valores cr\'iticos tabulados no son v\'alidos. Una soluci\'on habitual es
restar el \textbf{common time effect / efecto temporal com\'un}:

\begin{equation}
  \bar{y}_t = \frac{1}{n}\sum_{i=1}^{n} y_{it}
\end{equation}

y estimar las regresiones usando $y_{it}^* = y_{it} - \bar{y}_t$.

\subsection{Limitaciones}

\begin{enumerate}
  \item El rechazo de $H_0$ no identifica cu\'ales series son estacionarias.
  \item Los valores cr\'iticos son sensibles a c\'omo se toma el l\'imite
        asint\'otico ($n \to \infty$, $T \to \infty$, o ambos simult\'aneamente).
  \item La correlaci\'on contempor\'anea de los residuos invalida los
        valores cr\'iticos est\'andar.
\end{enumerate}

\subsection{Otras Pruebas de Panel}

La prueba de \textbf{Maddala y Wu (1999)} es similar a la IPS pero requiere
valores cr\'iticos generados por \textit{bootstrap}. La prueba de
\textbf{Levin, Lin y Chu (2002)} impone la hip\'otesis alternativa m\'as
restrictiva $\gamma_1 = \gamma_2 = \cdots = \gamma_n$.

%%% =========================================================
\section{Tendencias y Descomposiciones Univariadas}
%%% =========================================================

Los hallazgos de Nelson y Plosser plantean la pregunta de c\'omo descomponer
una serie DS en sus componentes permanente (tendencia) y transitorio (ciclo).

\subsection{Descomposici\'on de Beveridge-Nelson (BN, 1981)}

Beveridge y Nelson (1981) demuestran que cualquier proceso ARIMA$(p,1,q)$
puede descomponerse en la suma de un paseo aleatorio con deriva y un
componente estacionario. La \textbf{tendencia estoc\'astica} $\mu_t$ se
define como el valor esperado condicional del l\'imite de la funci\'on de
pron\'ostico de largo plazo:

\begin{equation}
  \mu_t = \lim_{s\to\infty}\bigl(E_t\,y_{t+s} - a_0\,s\bigr)
\end{equation}

Para el caso ilustrativo ARIMA$(0,1,2)$, el componente estacionario
(irregular) es:

\begin{equation}
  y_t - \mu_t = -(\beta_1+\beta_2)\,\varepsilon_t - \beta_2\,\varepsilon_{t-1}
\end{equation}

y la tendencia evoluciona como un paseo aleatorio:

\begin{equation}
  \Delta\mu_t = a_0 + (1+\beta_1+\beta_2)\,\varepsilon_t
\end{equation}

Una propiedad caracter\'istica de la descomposici\'on BN es que la
correlaci\'on entre las innovaciones de la tendencia y del componente
irregular es siempre igual a $-1$.

El procedimiento operativo comprende tres pasos:
\begin{enumerate}
  \item Identificar y estimar el modelo ARMA para $\Delta y_t$
        (metodolog\'ia Box-Jenkins).
  \item Para cada $t$, calcular los pron\'osticos $s$-pasos adelante y
        construir la tendencia: $\mu_t = y_t + \sum_{s=1}^{\infty}E_t\,\Delta y_{t+s}$.
  \item El componente irregular se obtiene como $y_t - \mu_t$.
\end{enumerate}

\subsection{Descomposici\'on de Hodrick-Prescott (HP, 1997)}

Hodrick y Prescott (1997) extraen la tendencia $\{\mu_t\}$ minimizando
la suma de cuadrados ponderada:

\begin{equation}
  \frac{1}{T}\sum_{t=1}^{T}(y_t-\mu_t)^2
  +\frac{2\lambda}{T}\sum_{t=2}^{T-1}
   \bigl[(\mu_{t+1}-\mu_t)-(\mu_t-\mu_{t-1})\bigr]^2
\end{equation}

donde $\lambda$ es el \textbf{smoothing parameter / par\'ametro de
suavizamiento}:

\begin{itemize}
  \item $\lambda = 0$: la tendencia coincide con la serie observada
        ($\mu_t = y_t$).
  \item $\lambda \to \infty$: la tendencia converge a una funci\'on lineal
        del tiempo.
  \item $\lambda = 1600$ es el valor convencional para datos trimestrales.
\end{itemize}

El filtro HP puede aplicarse simult\'aneamente a m\'ultiples variables,
garantizando que todas compartan la misma forma de tendencia. Sin embargo,
se ha demostrado que introduce fluctuaciones espurias en el componente
c\'iclico; es m\'as apropiado cuando $\{y_t\} \sim I(2)$.

%%% =========================================================
\section{Conclusiones}
%%% =========================================================

El cap\'itulo establece la distinci\'on fundamental entre tendencias
determin\'isticas y estoc\'asticas y sus consecuencias para el an\'alisis
de series de tiempo:

\begin{itemize}
  \item La diferenciaci\'on elimina tendencias estoc\'asticas; la
        desestacionalizaci\'on determin\'istica elimina tendencias
        determin\'isticas. Aplicar el tratamiento err\'oneo genera
        problemas adicionales.
  \item Los estad\'isticos $t$ y $F$ convencionales son inv\'alidos para
        contrastar ra\'ices unitarias. Las pruebas DF proveen los
        estad\'isticos apropiados con distribuciones no est\'andar
        obtenidas por simulaci\'on de Monte Carlo.
  \item Las pruebas DF tienen potencia muy baja frente a alternativas con
        ra\'ices casi unitarias. La prueba DF-GLS (ERS) mejora
        sustancialmente la potencia al estimar la tendencia de forma
        m\'as eficiente.
  \item Los quiebres estructurales sesgan las pruebas DF hacia la no
        rejecci\'on; la prueba de Perron (1989) permite incorporar un
        quiebre conocido en el an\'alisis.
  \item Las pruebas de ra\'iz unitaria en panel combinan informaci\'on de
        $n$ series para aumentar la potencia, aunque el rechazo es menos
        espec\'ifico que en el caso univariado.
  \item La descomposici\'on BN impone una correlaci\'on de $-1$ entre
        innovaciones de tendencia e irregular. El filtro HP produce una
        tendencia m\'as suave pero puede introducir fluctuaciones espurias
        en el componente c\'iclico.
\end{itemize}

\end{document}
